Procedemos la demostración por \textbf{inducción sobre} $t$.

\bigskip

\boxed{\textbf{Caso base}} ($t = 1$):\\
Por definición de la inicialización del algoritmo, se tiene:
\[
\delta_j(1) = \pi_j \cdot B_{j,\,x^{(1)}}
\]

Esto es exactamente la probabilidad de comenzar en el estado $s_j$ y emitir la primera observación $x^{(1)}$ pues,
al no existir estados previos, el máximo sobre una secuencia vacía es simplemente la probabilidad conjunta del
primer estado y la primera observación:
\[
\delta_j(1) = p\!\left(x^{(1)},\; y^{(1)} = s_j\right)
\]
  
No hay estados anteriores que maximizar, por lo que el caso base se cumple trivialmente para $t=1$. $\checkmark$

\bigskip

\boxed{\textbf{Hipótesis de inducción}}:\\
Supongamos que para $t-1$, la variable $\delta_i(t-1)$ representa correctamente la probabilidad máxima para
cualquier estado $s_i \in S$:
\[
\delta_i(t-1) = \max_{y^{(1)} \cdots y^{(t-2)}} p\!\left(x^{(1)} \cdots x^{(t-1)},\ y^{(1)} \cdots y^{(t-2)},\ y^{(t-1)} = s_i\right)
\]

\bigskip

\boxed{\textbf{Paso inductivo}}:\\
Queremos demostrar que la recurrencia del algoritmo:
\[
\delta_j(t) = \max_{1 \leq i \leq N} \bigl[\delta_i(t-1)\cdot A_{ij}\bigr] \cdot B_{j,\,x^{(t)}}
\]
calcula correctamente el máximo hasta el tiempo $t$.\\

Desarrollamos directamente el máximo que queremos calcular:
\[
\max_{y^{(1)} \cdots y^{(t-1)}} p\!\left(x^{(1)} \cdots x^{(t)},\; y^{(1)} \cdots y^{(t-1)},\; y^{(t)} = s_j\right)
\]

Por la \textbf{independencia de emisión} del HMM, el factor $p(x^{(t)} \mid y^{(t)} = s_j) = B_{j,\,x^{(t)}}$ no
depende de los estados anteriores y puede extraerse del máximo:
\[
=\; B_{j,\,x^{(t)}} \cdot \max_{y^{(t-1)}}
\left[ A_{y^{(t-1)},\,j} \cdot \max_{y^{(1)} \cdots y^{(t-2)}} p\!\left(x^{(1)}
  \cdots x^{(t-1)},\; y^{(1)} \cdots y^{(t-2)},\; y^{(t-1)}\right) \right]
\]

Por la \textbf{propiedad de Márkov}, la transición $p(y^{(t)}=s_j \mid y^{(t-1)}) = A_{y^{(t-1)},\,j}$ también puede
separarse del máximo interior. Aplicando la \textbf{hipótesis inductiva}, el término interior es exactamente
$\delta_{y^{(t-1)}}(t-1)$:
\[
=\; B_{j,\,x^{(t)}} \cdot \max_{1 \leq i \leq N} \left[ \delta_i(t-1)\cdot A_{ij} \right]
\]
Esto es exactamente la recurrencia de Viterbi, por lo que la igualdad se cumple para $t$. $\checkmark$

\bigskip

\boxed{\textbf{Optimalidad de la cadena completa y backtracking}}:\\
Por el \textbf{Principio de Optimalidad de Bellman}:\\
Si una secuencia $y^{(1)}, \ldots, y^{(T)}$ es globalmente óptima, entonces su prefijo $y^{(1)}, \ldots, y^{(t)}$ es
óptimo para el subproblema hasta $t$, esto garantiza que $\delta_j(t)$ preserva el máximo correcto en cada paso
$t$.\\

Al finalizar la recursión, el estado óptimo final se identifica como:
\[
\hat{y}^{(T)} = \arg\max_{1 \leq j \leq N}\,\delta_j(T)
\]

Los apuntadores $\phi_j(t) = \arg\max_{i}\bigl[\delta_i(t-1)\cdot A_{ij}\bigr)$, almacenados durante la recursión,
reconstruyen la cadena completa hacia atrás. Por la inducción demostrada, esa cadena tiene probabilidad máxima
global:
\[
\hat{y}^{(1)}, \hat{y}^{(2)}, \ldots, \hat{y}^{(T)}
= \arg\max_{y^{(1)} \cdots y^{(T)}} p\!\left(x^{(1)} \cdots x^{(T)},\; y^{(1)} \cdots y^{(T)}\right)
\]

\bigskip
\bigskip

En la demostración es una aprovechamos dos ideas clave:
\begin{enumerate}
\item \textbf{La estructura del HMM permite separar el presente del pasado.}
  Gracias a la propiedad de Márkov y a la independencia de emisión, la probabilidad en el tiempo $t$ se puede
  escribir como:
  \[
  \underbrace{\delta_i(t-1)}_{\text{mejor camino hasta } t-1}
  \times
  \underbrace{A_{ij}}_{\text{transición}}
  \times
  \underbrace{B_{j,\,x^{(t)}}}_{\text{emisión}}
  \]
  El pasado queda completamente \textbf{resumido} en el escalar $\delta_i(t-1)$, no es necesario recordar cómo
  se llegó a ese valor, solo su resultado.

\item \textbf{El Principio de Optimalidad de Bellman.}
  Para llegar de forma óptima a un estado $s_j$ en el tiempo $t$, necesariamente se debió llegar de forma óptima
  al estado $s_i$ en $t-1$. En lugar de explorar todos los $|S|^T$ caminos posibles, basta con guardar el mejor
  valor acumulado hasta cada estado en cada instante, reduciendo la complejidad de
  $\mathcal{O}(|S|^T)$ a $\mathcal{O}(T \cdot |S|^2)$.
\end{enumerate}

El algoritmo de Viterbi puede verse como la búsqueda del camino más probable en un grafo de Trellis, donde en
cada columna temporal solo se conserva el mejor camino hacia cada nodo y se descartan todos los demás.
Los apuntadores $\phi$ son el mecanismo que permite recuperar ese camino óptimo al final, recorriéndolo
hacia atrás desde el estado final con mayor $\delta$.
